\section{On the wire}
\label{sec:tech:net}

Everything to this point is a property of the construction. This section is
about demonstrating it, because an argument that message sizes are
input-independent and a measurement that they are input-independent are
different claims, and only the second survives contact with an implementation.

\subsection{Three processes, not three objects}
\label{sec:tech:net:processes}

The servers are three separate operating-system processes. Each loads only
public data --- the item matrix $\itemembed$ and the catalogue --- and receives
one replicated share, which is useless alone. They never talk to each other:
in this slice every server-side operation is local, because $\itemembed$ is
public and $\mathbf{a}\cdot\itemembed$ is therefore a linear map on shares. The
only links are client-to-server, and that is a property of the protocol rather
than a simplification of the implementation.

Because a server serves one client to completion, nothing here is concurrent,
and \emph{no threading appears anywhere in this project}. Choosing separate
processes over threads made that true rather than merely likely.

\paragraph{Framing.} \textsc{Requirements} \S5 fixes the transport as ``plain
TCP, length-prefixed frames''. A frame is a four-byte little-endian length,
written with explicit shifts rather than a \texttt{memcpy} of a
\texttt{uint32\_t}, then exactly that many bytes.

Two details are load-bearing and both are tested. TCP is a byte \emph{stream},
so a 227-byte key handed to \texttt{send()} may arrive as 200 bytes and then
27; anything that assumes one \texttt{send()} equals one \texttt{recv()} works
on loopback with small payloads and fails on a real link, which is the worst
kind of bug because the demo passes on the development machine. The test
therefore \emph{forces} a split, by giving the listener a deliberately tiny
receive buffer, rather than hoping to observe one.

The second is that the declared length is untrusted input. It arrives over a
socket, so a peer chooses it, and reading it and then allocating that many
bytes is a four-byte denial of service. It is validated against a cap
\emph{before} it sizes anything --- the same discipline
\texttt{DpfKey::Deserialize} already applies, and the reason this parser is
covered by the hardened build alongside the other two.

\subsection{The distinguisher}
\label{sec:tech:net:distinguisher}

The phase's exit criterion originally asked for \texttt{tcpdump} on the server
links showing ``nothing but pseudorandom bytes''. That was amended, partly
because there is no \texttt{tcpdump} on the development platform, but mainly
because \emph{bytes looking random is not evidence}. Any fixed byte string
looks random if you do not know what to compare it against. The claim that
matters is that an adversary given the wire bytes cannot say which record was
fetched --- and the way to establish that is to build the adversary and measure
it losing.

\texttt{probe} records many independent queries for each of several fixed
record indices, calling $\Gen$ afresh every time, so the keys are independently
sampled; an adversary fed one key repeated many times would be measuring
nothing. Two claims are then tested separately, because they can fail
separately.

\paragraph{Length.} Every frame must be the same size whatever the index. This
follows from the construction --- key size depends only on the domain width ---
but a framing change could break it silently, so it is checked in the data.
Over \textbf{1800 recorded queries across three record indices, every frame is
228 bytes}.

\paragraph{Content.} A logistic regression is trained on the raw wire bytes and
scored on held-out data.

\begin{table}[h]
\centering
\small
\caption{Distinguishing between record indices from the wire bytes alone.
Chance is $0.5$, and it lies inside every interval. The shuffled-label control
runs the identical pipeline on deliberately destroyed labels: it agrees, which
is what shows the harness is measuring what it claims to.}
\label{tab:distinguisher}
\begin{tabular}{@{}lcc@{}}
\toprule
adversary & accuracy & 95\% CI \\
\midrule
index 0 vs 1234    & 0.490 & $[0.445,\ 0.534]$ \\
index 0 vs 777     & 0.494 & $[0.449,\ 0.538]$ \\
index 1234 vs 777  & 0.500 & $[0.455,\ 0.545]$ \\
\midrule
shuffled control   & 0.483 & $[0.439,\ 0.528]$ \\
\bottomrule
\end{tabular}
\end{table}

\begin{figure}[h]
\centering
\includegraphics[width=0.72\linewidth]{figures/fig_distinguisher}
\caption{The adversary failing. What matters is not that the points sit near
$0.5$ but that the chance line falls inside every interval; reporting a bare
``we got 49\%'' would be claiming more precision than 480 held-out trials
support.}
\label{fig:distinguisher}
\end{figure}

We state the reading carefully. This is not a proof, and a null result is not
the same as a security proof: it says a particular adversary, given a
particular feature representation, did not beat chance on this data. What it
rules out is the failure mode that actually occurs in practice --- a length or
byte-level pattern correlating with the query --- and it does so by
demonstrating an adversary losing rather than by asserting that bytes look
random.

Two things are still explicitly \emph{not} claimed. Inter-frame \emph{timing}
is not analysed; only sizes and contents are. And an observer who sees both
server links can reconstruct the record by design, which is what a $(2,2)$
scheme means and is not a weakness peculiar to this implementation.

\subsection{What it cost to say this}
\label{sec:tech:net:cost}

The networked path returns the same ten titles, in the same order, as the
in-process path, with every record still byte-exact --- which is the assertion
the demo makes rather than a claim we make about it. A full run of the
protocol, including uploading the shares and the mask and fetching ten records,
moves about 86\,KB up and 46\,KB down across the three links, framing included.

Almost all of that is the seen-item mask, which is a shared vector of length
$\items$ and therefore dwarfs the ten 227-byte keys. That is worth noticing
rather than hiding: it is an argument for masking client-side in this slice,
where it is anyway equivalent, and it is the first place a wide-area deployment
would look for a saving.
